\documentclass[11pt]{article}
\usepackage[utf8]{inputenc}
\usepackage[T1]{fontenc}
\usepackage{lmodern}
\usepackage{amsmath,amssymb,amsthm}
\usepackage{graphicx}
\usepackage{booktabs}
\usepackage{enumitem}
\usepackage{hyperref}
\usepackage[margin=1in]{geometry}
\usepackage{listings}
\usepackage{xcolor}
\lstset{basicstyle=\ttfamily\small,breaklines=true,frame=single,
        backgroundcolor=\color{gray!8}}
% Markdown-math compatibility: these come from this user's markdown convention
% (bare <,> break PyCharm's math preview, so \lt/\gt are used in source).
\providecommand{\lt}{<}
\providecommand{\gt}{>}

\title{Alexander and Thurston norms for word cocycles\\ \large a formal account}
\author{}
\date{}

\begin{document}
\maketitle

A Bourbaki-style (definition--theorem--proof) formalization of
\texttt{[[thurston-alexander-norm]]}. The algebraic backbone (the representation
$\mu$, the cocycle $S$, the coboundary $D$, the monodromy polynomial) is recalled
from the certified claims \textbf{C1--C11} of \texttt{structure.md}; the norms and
the inequality of \S4--\S5 are new and proved here. The inequality is also guarded by
a numerical battery, \texttt{thurston\_alexander\_check.py} (exits 0).

Throughout, all polynomials live over the integral domain $R=\mathbb Z[X,Y]$ and we
write $\langle\phi,v\rangle$ for the pairing of a linear form $\phi$ with a vector $v$.

\hrulefill

\section{The algebraic structures}

\textbf{1.1 Definition (grading group, monomials).} Let $H=\mathbb Z^2$ with basis
dual to $(X,Y)$; a monomial $X^aY^b$ has \emph{multidegree} $\deg(X^aY^b)=(a,b)\in H$.
The \emph{Newton support} of $f=\sum_{(a,b)}c_{ab}X^aY^b\in R$ is
$\operatorname{supp}(f)=\{(a,b):c_{ab}\ne0\}$ and its \emph{Newton polytope} is
$N(f)=\operatorname{conv}\operatorname{supp}(f)\subset H_{\mathbb R}:=H\otimes\mathbb R$.
A \emph{direction} is a linear form $\phi\in H_{\mathbb R}^{*}$, written
$\phi=(\alpha,\beta)$, $\langle\phi,(a,b)\rangle=\alpha a+\beta b$.

\medskip
\textbf{1.2 Definition (gap words).} Let $A=\mathbb Z_{\ge1}$ (the \emph{gap alphabet})
and $W=A^{*}$ the free monoid of \emph{gap words} $w=(d_0,\dots,d_{o-1})$, with
\emph{length} $o(w)=o$ and \emph{weight} $E(w)=\sum_id_i$. The \emph{partial sums} are
$e_k=\sum_{i<k}d_i$ ($0\le k\le o-1$), so $e_0=0$, $e_{k+1}-e_k=d_k$.

\medskip
\textbf{1.3 Definition (the Magnus--Borel representation).} Let
$\mu:W\to GL_2\big(R[Y^{-1}]\big)$ be the monoid homomorphism determined on letters by
$\mu(d)=\big(\begin{smallmatrix}Y&0\\1&X^{d}\end{smallmatrix}\big)$. By \textbf{C1/C2},
\[
\mu(w)=M(w)=\begin{pmatrix}Y^{o}&0\\ S(w)&X^{E}\end{pmatrix},\qquad
S(w)=\sum_{k=0}^{o-1}X^{e_k}Y^{\,o-1-k}.
\]
The image lies in the Borel (lower-triangular) subgroup; the diagonal carries the two
characters $w\mapsto Y^{o}$ and $w\mapsto X^{E}$, and the off-diagonal entry is the
\emph{cocycle} $S$.

\medskip
\textbf{1.4 Definition (cocycle, coboundary).} Let $\chi(w)=X^{E}/Y^{o}$ and
$\sigma(w)=S(w)/Y^{o}$. By \textbf{C5}, $\sigma(uv)=\sigma(u)+\chi(u)\sigma(v)$:
$\sigma$ is a crossed homomorphism, $\sigma\in Z^{1}(W,R_\chi)$. By \textbf{C6} the
coboundary of a constant $c$ is $w\mapsto c\,D(w)$ with the \emph{augmentation element}
\[
D(w)=X^{E}-Y^{o},
\]
so $[\sigma]\in H^{1}(W,R_\chi)$ is defined modulo $\langle D\rangle$. We call $D$ the
augmentation because $D(w)=\chi(w)Y^{o}-Y^{o}$ is the image of ``$w-1$'' under the
Borel character (Fox's fundamental formula); $S$ is the corresponding Fox-derivative
datum.

\hrulefill

\section{The Alexander module and the monodromy polynomial}

\textbf{2.1 Definition (Alexander module of a cyclic word).} Fix $w\in W$ and adjoin
the shift variable $t$, working over $\Lambda=R[Y^{-1}][t]$. The \emph{Alexander
module} of $w$ is
\[
\mathcal A_w=\operatorname{coker}\!\big(tI-M(w):\ \Lambda^{2}\to\Lambda^{2}\big).
\]

\medskip
\textbf{2.2 Proposition (monodromy polynomial = order ideal).} $\mathcal A_w$ is a
torsion $\Lambda$-module with order ideal generated by
\[
\Delta_w(t)=\det\!\big(tI-M(w)\big)=(t-X^{E})(t-Y^{o})=t^{2}-(X^{E}+Y^{o})t+X^{E}Y^{o}.
\]
\emph{Proof.} The order ideal of a module presented by a square matrix $P$ over a
domain is $(\det P)$. Here
$P=tI-M(w)=\big(\begin{smallmatrix}t-Y^{o}&0\\-S&t-X^{E}\end{smallmatrix}\big)$ is
lower-triangular, so $\det P=(t-Y^{o})(t-X^{E})$. This is exactly \textbf{C11}.
\hfill$\blacksquare$

\medskip
\textbf{2.3 Remark (fibered class).} Rotation of a cyclic word is conjugation by
$\mu(u)$ (\textbf{C3}), so $M(w)$ is the monodromy of the mapping torus of the shift,
and $\Delta_w$ is its \emph{fibered} Alexander polynomial. The eigenvalues are
$\lambda_1=Y^{o}$ (eigenvector $(D,-S)^{\top}$, \textbf{C4}) and $\lambda_2=X^{E}$
(eigenvector $(0,1)^{\top}$), whence $D=\lambda_2-\lambda_1$.

\hrulefill

\section{Newton-polytope seminorms}

\textbf{3.1 Definition.} For a nonempty finite $P\subset H_{\mathbb R}$ put
$\operatorname{w}_\phi(P)=\max_{p\in P}\langle\phi,p\rangle-\min_{p\in P}\langle\phi,p\rangle$,
the \emph{$\phi$-width} of $P$. For $0\ne f\in R$ set
$\operatorname{w}_\phi(f)=\operatorname{w}_\phi(N(f))$.

\medskip
\textbf{3.2 Proposition (seminorm).} For fixed $P$, $\phi\mapsto\operatorname{w}_\phi(P)$
is a seminorm on $H_{\mathbb R}^{*}$.
\emph{Proof.} $\operatorname{w}_\phi(P)=h_{P}(\phi)+h_{P}(-\phi)$ where
$h_P(\phi)=\max_{p\in P}\langle\phi,p\rangle$ is the support function of
$\operatorname{conv}P$. Support functions are sublinear (positively homogeneous and
subadditive); a sum of sublinear functions invariant under $\phi\mapsto-\phi$ is a
seminorm. \hfill$\blacksquare$

\medskip
\textbf{3.3 Proposition (Minkowski additivity).} For $f,g\in R\setminus\{0\}$,
$N(fg)=N(f)+N(g)$ (Minkowski sum), hence
$\operatorname{w}_\phi(fg)=\operatorname{w}_\phi(f)+\operatorname{w}_\phi(g)$.
\emph{Proof.} Ostrowski's theorem: over an integral domain the Newton polytope of a
product is the Minkowski sum of the factors' Newton polytopes (the $\phi$-extreme
monomials multiply without cancellation because $R$ is a domain). Support functions are
additive under Minkowski sums, and so is width. \hfill$\blacksquare$

\hrulefill

\section{The Alexander norm of a word}

\textbf{4.1 Definition (Alexander norm).} The \emph{Alexander norm} of $w$ in direction
$\phi$ is
\[
\lVert\phi\rVert_{A}(w):=\operatorname{w}_\phi\big(D(w)\big)=\operatorname{w}_\phi(X^{E}-Y^{o}).
\]
(We use the augmentation element $D$ --- the order datum of the class $[\sigma]$ modulo
which $S$ is defined, \S1.4 --- as the word's Alexander polynomial.)

\medskip
\textbf{4.2 Proposition (closed form).} $N(D)$ is the segment $[(E,0),(0,o)]$ and
\[
\lVert\phi\rVert_{A}(w)=|\alpha E-\beta o|.
\]
In the \emph{Collatz direction} $\phi=(\log2,\log3)$,
$\lVert\phi\rVert_A(w)=|E\log2-o\log3|=\big|\log(2^{E}/3^{o})\big|$.
\emph{Proof.} $D=X^{E}-Y^{o}$ has support $\{(E,0),(0,o)\}$; its $\phi$-width is
$|\langle\phi,(E,0)-(0,o)\rangle|=|\alpha E-\beta o|$. The specialization is immediate.
\hfill$\blacksquare$

\medskip
\textbf{4.3 Proposition (the multivariable Alexander polynomial is a parallelogram).}
In the variables $(X,Y,t)$ with direction $\phi=(\alpha,\beta,\gamma)$, $N(\Delta_w)$ is
the parallelogram with edge vectors $(E,0,-1)$ and $(0,o,-1)$, and
$\operatorname{w}_\phi(\Delta_w)=|\alpha E-\gamma|+|\beta o-\gamma|$.
\emph{Proof.} The four monomials of $\Delta_w$ have exponents
$(0,0,2),(E,0,1),(0,o,1),(E,o,0)$; their two diagonals share the midpoint
$(E/2,o/2,1)$, so they form a parallelogram with edges $(E,0,-1),(0,o,-1)$. By 3.3 (or
directly) its width is the sum of the two edge-widths $|\alpha E-\gamma|+|\beta o-\gamma|$.
Setting $\gamma=0$ recovers the factor structure
$\Delta_w|_{t=\text{const}}\sim (X^E)(Y^o)$ and the $D$-segment. \hfill$\blacksquare$

\hrulefill

\section{The Thurston norm of a word}

\textbf{5.1 Definition (Magnus path; per-letter increments).} The \emph{Magnus path} of
$w$ is the sequence of support points of $S$, $p_k=(e_k,\,o-1-k)\in H$ for
$0\le k\le o-1$. Its \emph{increments} are $\delta_k:=p_{k+1}-p_k=(d_k,-1)$,
$0\le k\le o-2$ (from $e_{k+1}-e_k=d_k$).

\medskip
\textbf{5.2 Definition (Thurston norm).} The \emph{Thurston norm} of $w$ in direction
$\phi$ is the total $\phi$-variation accumulated over the letters,
\[
\lVert\phi\rVert_{T}(w):=\sum_{k=0}^{o-1}\big|\langle\phi,(d_k,-1)\rangle\big|
=\sum_{k=0}^{o-1}|\alpha d_k-\beta|.
\]
Geometrically $\lVert\phi\rVert_T(w)=h_Z(\phi)$, the support function of the
\emph{letter zonotope} $Z=\sum_{k}[-\delta_k,\delta_k]$; in particular it is a seminorm
in $\phi$ (3.2 applied to each segment).

\medskip
\textbf{5.3 Remark (why ``Thurston'').} This is the complexity of the cyclic word as a
\emph{singular dual chain}: each letter contributes its $\phi$-defect
$|\alpha d_k-\beta|$, the discrete analogue of $\chi_-$ summed over the pieces of a
surface dual to $\phi$. The \emph{genuine} Thurston seminorm is the infimum over
representatives of the class, $\lVert\phi\rVert_T^{\min}=\inf_{w'\sim w}\lVert\phi\rVert_T(w')$;
Theorem 6.1 shows it equals $\lVert\phi\rVert_A$ (achieved by a $\phi$-monotone
reordering), matching McMullen's equality on fibered faces.

\hrulefill

\section{The Alexander--Thurston inequality}

\textbf{6.1 Theorem.} For every gap word $w$ and every direction $\phi$,
\[
\lVert\phi\rVert_{A}(w)\ \le\ \lVert\phi\rVert_{T}(w),
\]
with equality if and only if the per-letter defects $\alpha d_k-\beta$
($0\le k\le o-1$) all have the same sign (equivalently, the partial sums
$k\mapsto\langle\phi,p_k\rangle$ are monotone).

\medskip
\emph{Proof.} Write $x_k=\alpha d_k-\beta$. Since $\sum_{k=0}^{o-1}d_k=E$ and the sum
has $o$ terms,
\[
\sum_{k=0}^{o-1}x_k=\alpha E-\beta o=\langle\phi,(E,0)-(0,o)\rangle,
\]
so by 4.2, $\lVert\phi\rVert_A(w)=|\alpha E-\beta o|=\big|\sum_{k}x_k\big|$. The
triangle inequality gives $\big|\sum_k x_k\big|\le\sum_k|x_k|=\lVert\phi\rVert_T(w)$.
Equality in the triangle inequality holds iff all $x_k$ share a sign. \hfill$\blacksquare$

\medskip
\textbf{6.2 Geometric restatement.} $\lVert\phi\rVert_A\le\lVert\phi\rVert_T$ for all
$\phi$ is equivalent to the inclusion of bodies $N(D)-N(D)\subseteq Z$ (difference body
of the Alexander segment inside the letter zonotope), since the two sides are the
support functions of these bodies. The inclusion holds because the segment's difference
vector is $\sum_k\delta_k\in\sum_k[-\delta_k,\delta_k]=Z$.

\medskip
\textbf{6.3 Proposition (cocycle refinement).} The cocycle's own spread is also
controlled:
\[
\operatorname{w}_\phi\big(S(w)\big)\ \le\ \sum_{k=0}^{o-2}|\alpha d_k-\beta|\ \le\ \lVert\phi\rVert_T(w).
\]
\emph{Proof.} With $\psi_k=\langle\phi,p_k\rangle$,
$\operatorname{w}_\phi(S)=\max_k\psi_k-\min_k\psi_k$. If the max/min occur at indices
$M,m$ then
\[
|\psi_M-\psi_m|=\Big|\sum_{k=m}^{M-1}(\psi_{k+1}-\psi_k)\Big|
\le\sum_{k=0}^{o-2}|\psi_{k+1}-\psi_k|=\sum_{k=0}^{o-2}|\alpha d_k-\beta|,
\]
the first $o-1$ terms of $\lVert\phi\rVert_T$. \hfill$\blacksquare$

\medskip
\textbf{6.4 Corollary (genuine Thurston = Alexander; fibered equality).}
$\lVert\phi\rVert_T^{\min}(w)=\lVert\phi\rVert_A(w)$: reorder $w$ so that all letters
$\ge\beta/\alpha$ precede those $<\beta/\alpha$; the reordered word $w'$ (same multiset
of letters, same $E,o$, hence same $\lVert\cdot\rVert_A$) is $\phi$-monotone, so by 6.1
$\lVert\phi\rVert_T(w')=\lVert\phi\rVert_A(w')=\lVert\phi\rVert_A(w)$, and no
representative can go below $\lVert\phi\rVert_A$ by the theorem. This is the
word-analogue of McMullen's equality $\lVert\phi\rVert_A=\lVert\phi\rVert_T$ on a
fibered face. \hfill$\blacksquare$

\hrulefill

\section{Interpretation, scope, and what is open}

\begin{itemize}[leftmargin=*]
  \item \textbf{The Collatz reading.} $\lVert\phi\rVert_A(w)=|\log(2^{E}/3^{o})|$ is
  the \emph{log size barrier} ($|D|=|2^{E}-3^{o}|$, the Ellison floor,
  \texttt{[[ellison1971-pillai]]}); $\lVert\phi\rVert_T(w)=\sum_k|\log(2^{d_k}/3)|$ is
  the \emph{total per-step imbalance}. Theorem 6.1 reads: the net imbalance never
  exceeds the total --- a cycle's $D$ is small only through cancellation of per-letter
  imbalances, never because the letters are individually balanced.
  \item \textbf{Equality locus.} Equality holds iff all gaps lie on one side of
  $\beta/\alpha=\log_2 3\approx1.585$, i.e. all $d_k=1$ or all $d_k\ge2$. The all-equal
  circuits $d^{o}$ (the $C=0$ stratum, \textbf{P1}) are the archetype; they are the
  $\phi$-monotone (geodesic) representatives where Alexander $=$ Thurston.
  \item \textbf{Span vs. height (do not conflate).} Both norms here are \emph{span}
  (Newton-width) data; they are unrelated to the \emph{height} (Mahler measure $=$
  entropy) complexity of \texttt{[[cocycle-complexity]]}. A cyclotomic factor inflates a
  height-complexity of $1$ while contributing fully to the width. The two axes meet only
  on the fibered cone (Fried/McMullen), which for words is 6.4.
  \item \textbf{Status.} \S1--\S2 are certified (\textbf{C1--C11}); 3.2--3.3 are
  standard (support functions, Ostrowski); 4.2, 4.3, 5.2, and Theorem 6.1 with 6.3--6.4
  are proved above and battery-checked
  (\texttt{thurston\_\allowbreak alexander\_\allowbreak check.py}, 13728 cases, exit 0).
  \item \textbf{Open.} (i) Promote $\lVert\cdot\rVert_A$ to the full multivariable
  Alexander polynomial $\Delta_w$ (4.3) and prove a $3$-variable inequality against a
  Turaev $2$-complex norm of the word's presentation. (ii) Extend to the \emph{core}
  $C(u,v)$: does $\operatorname{w}_\phi(C)$ obey an analogous bound, and where do its
  width and Mahler measure disagree (the place to hunt, per
  \texttt{[[thurston-alexander-norm]]} \S5)? (iii) Identify the genuine Thurston
  seminorm with a $\min$-$\chi_-$ over an explicit cell structure.
\end{itemize}

\end{document}
